\documentclass[11pt]{article}
\usepackage{amsmath,amssymb,amsthm,geometry}\geometry{margin=1in}
\newtheorem{definition}{Definition}\newtheorem{theorem}[definition]{Theorem}\newtheorem{proposition}[definition]{Proposition}
\title{Boundary-row algebras: a non-rigid native feedback interface}
\date{16 November 2026}
\begin{document}\maketitle
\paragraph{Audit.} Standalone presentation of the command-indexed native theorem. The original version SHA256 \texttt{38cb2f2001495d92ed3c8c330dff6d7892589b787aef32b7c6621dd44e1f9cab} received independent conditional review \texttt{246224c1f59e72864f83e40eb46d52787c55329c0f4f96b46b124d5e6eaafdfb}; this version incorporates its explicit interface pins. The construction below is an independently defined category of opaque operational systems, not the finite-word category. It gives a native initiality theorem for \emph{command-indexed full writer families}. It does not yet give initiality for a doctrine requiring one shared writer-library map across aliased schemas. That distinction is an explicit kernel, not a claim that unused coefficients may be dropped. All shapes and arrows of each command's full family are retained. Scope writer and terminal polynomial grades remain separate. Port-free deterministic seeds are imposed; replacing this with finite silent initialization requires another proof.

\paragraph{Pinned dependencies.} The sources read are World-Rebased Multi-Active Open Assembly v3, SHA256 \texttt{bbc0006ed160be930687306e2f57b6ab57d111e93b13e6c92b4f80dc881d60cd}; Whole-Scope Block Decomposition v3, \texttt{d3a9f7dfb58de84e685ec41983b06ed4914fa615d76b6aa23040627e328b1fb3}; Feedback Shared Rebase Research Notes v3, \texttt{71891637a7c10a0618146797fd4e14129de68e32426287e033ec607c88e10416}; and Pure Wires and Relative Shared-World Assembly. Only the finite closure, whole-law block decomposition, and actual limit comparisons are used, not an unproved opaque-comparison or feedback theorem.

\section{An independent category}
Fix worlds $\mathcal W$, recursive input colors $U$, external output colors $B$, and a boundary library. A boundary schema $e$ at $w$ consists of a finite carrier $D_e$, an actual arrow $d:w\to w_d$ for every $d\in D_e$, context objects $L_{e,d}$, a full writer family $H_e$, and a body $\beta_e:D_e\to |H_e|$. Distinct boundary witnesses can have the same world arrow. No boundary schema is silently refined. Fix categories of augmented context coefficients and full writer families in which the finite slice products below exist and are chosen functorially (or retain their canonical coherent isomorphisms).

An object $X$ has arbitrary sets $S_X(w,u)$, distinguished seeds, functorial rebase, stable finite ordered external ports, and completion markers stable under rebase. Its independently full terminal collection $C_X$ supplies a terminal body on every completed state. Passive rebase preserves the completed terminal body, using the literal stable port identification. This collection is a colored weighted polynomial, with actual maps on every coefficient, including unrealized shapes. Completed states have no active commands. Seeds have no ports; they need not be natural under rebase.

For each state $s$ and schema $e$, its active menu is a set. A command $c$ contains a finite normalized row $(Z_c,q_c)$, a boundary map $b_c:Z_c\to D_e$, actual context coefficients $K_{c,z}\to L_{e,b_cz}$, a whole writer family $Q_c\to H_e$, and a body $a_c:Z_c\to |Q_c|$ over $\beta_e b_c$. It has a successor at world $w_{b_cz}$, preserving every old external port. No participant decomposition of this command is required. In particular an opaque command may implement a genuinely simultaneous network event.

The passive response at a specified boundary witness $d$ is $R_d(s)$ with context coefficient $L_{e,d}$ and writer body $\beta_e(d)$ in the whole family $H_e$. It is not a member of the active menu and carries no paid step. Its distribution, when packaged across $D_e$, is imposed by the surrounding joint, never independently sampled.

A morphism $f:X\to Y$ has a state function in every world/input fiber, preserves seeds, rebase and completion in both directions, supplies order- and label-preserving external-port bijections, and has one full terminal collection map. These bijections commute with passive rebase and with old-port inclusions under every active raw outcome, including zero atoms. The terminal body square uses these same bijections. At each state and active command it chooses a target command over the same $e$, a boundary-preserving raw function with exact pushforward, actual context arrows over $L$, and an actual full-family arrow $Q_c\to Q_{f(c)}$ over $H_e$. The successor and full-body squares commute on all raw atoms. No equality is imposed merely because two command-indexed families happen to be isomorphic. Identities and composition are fieldwise. Passive responses are mapped by identity on the fixed boundary data and the rebase square.

This definition precedes the constructor. State maps need not preserve an invocation address, word length, participant count, or hidden implementation. The algebra structure below is an ordinary morphism in this category; its whole-command arrows are supplied data, not inferred from primitive marginals. Native maps may identify states or commands, but not omit source commands or raw zero atoms.

\section{The root endofunctor}
Fix a primitive $P$ with root states having finite ordered persistent slots, each either an external $B$ port or a recursive $U$ port. Every root state and raw successor has finitely many slots. Passive rebase preserves them. Root seeds have no slots. Root command data use the boundary interface above. Root completion supplies a whole colored collection $C_P$ that records both kinds of slots. Write $G_P(C)$ for its ordinary colored polynomial substitution with $C$ in recursive slots and identity external leaves; keep all formal shapes and coefficient tensors.

A state of $T_PX$ is a root state $p$ and, at each existing recursive port $j$ of $p$, a child $x_j$ of matching input at the same world. Its external ports are the planar root external ports and the child external ports, with tagged paths. Completion means root and every child complete. Rebase is componentwise. The seed is the root seed, with no children. The full terminal grade is $G_P(C_X)$.

A paid command at $(p,(x_j))$, over $e$, selects an active root command or a passive root response, and independently an active child command or passive child response for each child, with at least one active choice. Its raw carrier is the entire matching product of the selected active raw rows and passive $D_e$ coordinates over $D_e$. Its law is any normalized law on this carrier whose active marginals are their selected laws. The passive marginals are induced by that one law. Its context coefficient is the slice product of the selected $K$ objects and passive $L$ objects. Its whole writer family is the slice product of the selected whole $Q$ families and passive $H_e$. Its body is the induced full cone body. There is no independent draw attached to a passive block.

All old children receive their selected active successor or their passive rebase under the same boundary witness. An active root may create recursive ports; each new child is the seed of $X$ at that port's input and the actual post-world. New children do not join the already selected scope. Root outputs and old child outputs persist. This specifies the whole state, not just its external trace.

On a native map $f$, apply its state map to children. Transport selected child commands using its menu map, use the product raw function, and name the pushed joint law as the target command. Pushforward has the required target marginal laws. The actual context and whole-family maps are the induced slice-product arrows. The root and passive boundary factors are identities. The terminal grade map is $G_P(C_f)$.

\begin{proposition}
These formulas define an endofunctor on the category just specified, up to chosen canonical limit coherence. They do not require a primitive participant map for an opaque child command.
\end{proposition}
\begin{proof}
Every generated row is finite. Exact marginal preservation follows by composing pushforwards with projections. All boundary squares commute by the native map definition. Thus the actual coefficient and whole-family cone maps exist. The post-square for an old active child is its native successor square; for an old passive child it is rebase equivariance; for a new child it is preservation of seeds at the post-world. Port persistence is checked before applying the port bijections. Terminal body preservation is colored polynomial substitution of the full body maps. Identity and composition follow from pushforward functoriality and uniqueness of the actual limiting arrows. In particular a target menu containing all allowed couplings, rather than a fixed chooser, is essential to this proof.
\end{proof}

\section{Finite native carrier and fold}
Let $F_P$ have all finite compatible rooted $P$-invocation trees at a common world/input, not only reachable ones. Every recursive root port has exactly one child. Its seed is the one-vertex root seed. Rebase is nodewise, ports are planar external leaves, and completion means all nodes complete. Its paid commands choose a nonempty finite set of tree vertices and one primitive root command at each, all over the same $e$, together with an arbitrary law on the full matching raw carrier with those primitive marginals. Actual context coefficients and full command writer families are the flat slice products. Passive nodes contribute boundary responses, not additional events or samples. Births insert seed vertices at the actual post-world.

Its terminal grade $C_F$ is the collection of all finite formal colored $C_P$-trees, not only those realized by completed configurations. The coefficient of a shape is the ordered tensor of its vertex coefficients. This is the standard initial $G_P$-algebra: a shape map and actual coefficient map into any $G_P$-algebra are defined by finite-tree recursion; uniqueness follows by the same recursion and monoidal coherence. No joint writer coefficient is identified with that tensor.

Root splitting partitions each flat scope into the root and immediate child subtrees. Empty restrictions use the passive branch. Each nonempty child restriction has its uniquely induced marginal law. The retained outer coupling is the pushforward of the entire flat law, not its product of marginals. Conversely the selected child laws and outer coupling uniquely reconstruct the flat law. The pinned block theorem gives actual raw, context and full-family comparisons, including zero atoms and unused shapes. Together with grafting and the full terminal collection constructor these give a native isomorphism $\iota:T_PF_P\to F_P$.

\begin{theorem}[Command-indexed native initiality]
For every native algebra $\alpha:T_PX\to X$ in the preceding independently specified category, there is a unique native algebra morphism $f:F_P\to X$. This theorem concerns command-indexed full writer families and deterministic port-free seeds.
\end{theorem}
\begin{proof}
Define the state function by structural recursion,
\[
f(p,(t_j))=\alpha(p,(f(t_j))).
\]
Define the full terminal collection map by the independently full $G_P$-algebra recursion induced by the terminal-grade component of $\alpha$. Neither recursion deletes unreachable configurations or unrealized shapes.

For a source command at a finite tree, split the whole command at the root. Recursively map every nonempty child restriction; a passive child uses the canonical boundary response. Apply the actual product raw map and push the retained outer joint law. Assemble the actual context and whole-family maps by the limiting cones, and then compose with the corresponding whole-command fields of $\alpha$. This recursively defines the menu choice, raw function, context arrows and full-family map. The recursion decreases the size of the \emph{prescope} child tree, even when the command's successor contains more vertices. Whole-command fields of $\alpha$ determine opaque target responses; no primitive-projection uniqueness assumption is imposed on $\alpha$ itself.

Verify native laws by simultaneous induction on the source prescope tree. Rebase preserves its shape, so rebase equivariance follows from the child induction and $\alpha$'s native rebase square. Seed preservation holds because the seed has one port-free root and $\alpha$ preserves the constructor seed. Existing active child successor squares follow from the inductive native law for that smaller prescope tree, which quantifies over \emph{all} its raw outcomes and their possibly larger successors. Existing passive child squares follow from rebase equivariance. Fresh child squares use the already proved seed equation at the actual post-world. Applying $\alpha$'s successor square then proves the entire parent square. This does not induct on successor height or assume a step decreases height. Boundary, exact law and whole-body equations follow under composition and the actual slice-product assembler. Port/completion equations are preserved by $T_P$ and $\alpha$. Terminal bodies use the separate full-grade recursion.

The constructor equation $f\iota=\alpha T_Pf$ holds by these definitions on every native field, including all whole-family arrows. Conversely any algebra morphism satisfies this equation on states and terminal grades, forcing the two structural recursions. On a source command the same equation forces the root split, the recursively determined child command maps, the pushed full coupling, the induced actual slice-limit arrows, and finally the supplied fields of $\alpha$. Induction on the source tree forces its menu, raw, context and complete writer-family fields. Thus no two native algebra morphisms differ on an unused shape or hidden joint coefficient. This is equality of full native morphisms, not equality only of primitive restrictions or traces.
\end{proof}

\section{Meaning, obstruction boundary, and review obligations}
The algebra category is not a rigid occurrence category: $\alpha$ can absorb a root and any finite number of opaque child states into a single opaque state. Its command map can change internal participant organization because that organization is not an object field. The price is explicit whole-row and whole-writer arrows, which are exactly what primitive marginal maps failed to determine. A simple non-rigidity witness uses the trivial world, one recursive input color, no commands or external output ports, and unit terminal coefficients. A primitive has a completed nullary seed state and a completed unary-recursive state. All finite chains map to the single completed commandless port-free state of an opaque system whose entire terminal collection has one nullary shape with unit coefficient; the whole grade map sends every chain shape to that nullary shape using the canonical unit-coefficient arrow. The algebra absorbs any unary layer; no fixed-word-length category admits that state identification by componentwise maps alone. This witness establishes breadth, not a shared-world stress test; the theorem's shared scopes are handled by its full coupling construction.

\paragraph{Library alias obstruction.} A stronger source doctrine may identify a schema used by two different commands and require one writer-family map for both. The construction above gives a full arrow for each command, but does not prove these arrows equal for two alias occurrences. Equality on realized bodies is insufficient: even a singleton raw row can leave an unused coefficient direction unconstrained. A source promotion must either (i) supply a global schema functor on objects and native arrows and show the recursive scope assembler respects every alias relation, or (ii) retain free scope-plan schema syntax and prove a descent theorem for the intended quotient. Neither follows from limits at individual commands. This is the current exact residual, rather than an assertion of full source-doctrine initiality.

\paragraph{Other nontransporting changes.} Restricting the joint admission policy needs restriction and pushforward closure proofs. Permitting non-port-free seeds needs a finite initializer comparison and a seed-fold proof; no seed-naturality under world rebase is available. Identifying different boundaries by world endpoint loses the passive writer witness. Replacing actual products by arbitrary supplied cones loses the unique assembler. Merging the scope writer and terminal polynomial grades remains prohibited by the pinned joint-to-tensor obstruction. None of these changes is included in the theorem.

\paragraph{Independent review targets.} Check that the category's command-indexed full-family field matches the desired observer; check the all-config induction on successor squares; check that command normalization removes only passive coordinates forced by the common boundary and retains all active tuple zeros; and check the full terminal collection recursion in the exact source colored polynomial doctrine. The result is a substantive non-rigid algebra candidate, not automatically promoted into the shared-library source theorem.
\end{document}
